% !TeX root = ../main.tex
\renewcommand{\abstractname}{Abstract}
\begin{abstract}
As part of the transition towards Industry 4.0, industrial assets and autonomous guided vehicles have evolved into systems capable of transmitting continuous streams of real-time telemetry, generating a vast amount of information. However, simply collecting this data does not, in itself, provide any operational value. Only through appropriate processing and analysis is it possible to transform this data into useful knowledge for strategic decision-making.
This is where the idea for this project came from, developed in collaboration with EDAG, a leading organisation in the industrial sector.\\

Despite the progress made in predictive maintenance in recent years, there remains a significant gap between theoretical studies on anomaly detection and their effective application in production environments. This limitation is mainly due to the complexity of collecting high-quality data, the dynamic nature of industrial environments and the need to process information with low latency in order to take action before a failure occurs.\\

To address this issue, this project presents a comprehensive solution for real-time anomaly detection across a fleet of mobile robots. The system comprises a distributed infrastructure for continuous data ingestion and processing, a historical storage layer optimised for time series, and an analytics module based on machine learning for identifying anomalous behaviour. The system is rounded off by a visual monitoring and alerting tool for plant staff, developed using a clean, decoupled and scalable architecture.\\
\medskip
\noindent\textbf{Keywords:} IIoT, Predictive Maintenance, Anomaly Detection, Machine Learning, AGV, AMR, Industrial Robots, Industry 4.0, Data Lakehouse, Medallion Architecture.
\end{abstract}
